\begin{tikzpicture}[font=\small, node distance=12mm and 18mm]

% Column x positions (for neat horizontal layout)
\coordinate (col1) at (0,0);
\coordinate (col2) at (2.25,0);  % first backbone layer
\coordinate (col3) at (8.10,0);  % similarity module centers
\coordinate (col5) at (13.8,0);  % task-specific layers (right)

% Row y position (single task)
\coordinate (rowA) at (0,0);

% ---------- Column 1: input dataset ----------
\node[data] (dA) at ($(col1)+(rowA)$) {$X$};

% ---------- Column 2: first backbone layer ----------
\node[back, minimum width=2.5cm, minimum height=2.5cm,
      label={[CFrame]below: Input Module}] (bA) at ($(col2)+(rowA)$) {$\mathbf{f}^{(\mathrm{in})}_t$};
\draw[arrow] (dA) -- (bA);

% ---------- Similarity Module ----------
\node[sim, minimum width=2.5cm, minimum height=2.5cm,
      label={[CFrame]below: Similarity Module}] (simA) at ($(col3)+(rowA)$) {$\mathbf{f}^{(\mathrm{sim})}_t$};

% ---------- Input feature circles (OUTSIDE, to the LEFT of the module) ----------
\foreach \i/\lab in {1/$h_{t,1}^{(in)}$,2/$h_{t,2}^{(in)}$,3/$h_{t,3}^{(in)}$}{
  \node[fcin] (Ain\i) at ($(simA.west)+(-1.2,2.4-1.2*\i)$) {\scriptsize \lab};
}

% Connect backbone -> EACH input circle
\foreach \i in {1,2,3}{
  \draw[arrow] (bA.east) -- ($(Ain\i.west)+(-0.25,0)$) -- (Ain\i.west);
}

% Horizontal arrows from all input circles to the similarity module
\foreach \i in {1,2,3}{
  \draw[-{Latex[length=1.6mm]}, CFrame, thick] 
      (Ain\i.east) -- ($(simA.west)+(0,1.5-0.8*\i)$);
}

% ---------- Output feature circles (OUTSIDE, to the RIGHT of the module) ----------
\node[fcout] (Aout1) at ($(simA.east)+(1.2,0.7)$) {\scriptsize $h_{t,1}^{(sim)}$};
\node[fcout] (Aout2) at ($(simA.east)+(1.2,-0.7)$) {\scriptsize $h_{t,2}^{(sim)}$};

% Arrows from module to EACH output circle
\draw[-{Latex[length=1.6mm]}, CFrame] ($(simA.east)+(0,0.70)$) -- (Aout1.west);
\draw[-{Latex[length=1.6mm]}, CFrame] ($(simA.east)+(0,-0.70)$) -- (Aout2.west);

% ---------- Column 5: task-specific layers ----------
\node[back, minimum width=2.5cm, minimum height=2.5cm,
      label={[CFrame]below: Output Module}] (hA) at ($(col5)+(rowA)$) {$\mathbf{f}^{(\mathrm{out})}_t$};

% Connect EACH output circle -> task-specific layers
\draw[arrow] (Aout1.east) -- (hA.west);
\draw[arrow] (Aout2.east) -- (hA.west);

% ---------- ALE explainer box with top label (ABOVE THE FIGURE) ----------
\node[
  draw=CFrame,
  rounded corners=6pt,
  fill=white,
  align=left,
  inner sep=6pt,
  minimum width=3.0cm,
  minimum height=2.8cm,
  above=10mm of simA,
  label={[CFrame]above:\textbf{ALE for $(j_1 \rightarrow k_1)$}}
] (aleBox) {};

% ---- ALE axes and curve (clipped to stay inside the axes) ----
\path (aleBox.south west) ++(0.55,0.55) coordinate (ax0);
\def\aleW{2.4}   % x-axis length
\def\aleH{1.4}   % y-axis length

% Axes
\draw[->, thin] (ax0) -- ++(0,\aleH) node[above] {\scriptsize $h_{t,1}^{(sim)}$};
\draw[->, thin] (ax0) -- ++(\aleW,0) node[right] {\scriptsize $h_{t,1}^{(in)}$};

% Clip region = axes box
\begin{scope}
  \clip (ax0) rectangle ($(ax0)+(\aleW,\aleH)$);
  % ALE curve (clipped inside axes)
  \draw[CFrame, thick]
    ($(ax0)+(0.10,0.20)$)
    .. controls +(0.60,0.60) and +(-0.40,0.20) ..
    ($(ax0)+(1.20,1.00)$)
    .. controls +(0.40,-0.30) and +(0.15,0.35) ..
    ($(ax0)+(2.30,0.65)$);
\end{scope}

% Dashed pointers from selected pair (Ain1, Aout1) up to the ALE explainer
\draw[densely dashed, CFrame, -{Latex[length=1.2mm]}]
  (Ain1.north) |- ($(aleBox.west)$);
\draw[densely dashed, CFrame, -{Latex[length=1.2mm]}]
  (Aout1.north) |- ($(aleBox.east)$);

\end{tikzpicture}
